% !TEX program = xelatex
\documentclass[11pt,a4paper]{article}

\usepackage[UTF8]{ctex}
\usepackage[margin=2.2cm]{geometry}
\usepackage{amsmath, amssymb, amsthm}
\usepackage{listings}
\usepackage{xcolor}
\usepackage{tikz}
\usepackage{booktabs}
\usepackage{multirow}
\usepackage{longtable}
\usepackage{hyperref}
\usetikzlibrary{positioning, arrows.meta, shapes.geometric, fit, calc, decorations.pathreplacing}

\hypersetup{
  colorlinks=true,
  linkcolor=blue!60!black,
  urlcolor=teal!70!black,
}

\lstdefinelanguage{Rust}{
  morekeywords={fn,let,mut,pub,struct,impl,use,mod,self,Self,Option,Some,None,return,if,else,for,while,loop,match,break,continue,as,where,trait,type,enum,move,ref,const,static,unsafe,extern,crate,super,dyn,async,await,box,in},
  sensitive=true,
  morecomment=[l]{//},
  morecomment=[s]{/*}{*/},
  morestring=[b]",
}

\lstset{
  language=Rust,
  basicstyle=\ttfamily\footnotesize,
  keywordstyle=\color{blue!60!black}\bfseries,
  commentstyle=\color{green!40!black}\itshape,
  stringstyle=\color{red!60!black},
  numbers=left,
  numberstyle=\tiny\color{gray},
  numbersep=8pt,
  frame=single,
  framerule=0.4pt,
  rulecolor=\color{gray!40},
  breaklines=true,
  breakatwhitespace=true,
  showstringspaces=false,
  tabsize=4,
  columns=fullflexible,
  keepspaces=true,
}

\title{\textbf{halfedge\,: 内置属性模块设计文档} \\ \large \texttt{src/builtin\_attrs.rs}}
\author{halfedge 库}
\date{2026-07-05}

\begin{document}
\maketitle
\tableofcontents

\section{模块定位}

\texttt{builtin\_attrs.rs} 在 \texttt{property.rs} 泛型属性系统之上，
为\textbf{常见几何属性}（顶点法向 / UV / 颜色 / 面法向）提供强类型 Newtype 包装与
便捷函数，并实现\textbf{OBJ 属性感知 IO}（\texttt{v/vt/vn} 解析与序列化）。

\section{为什么需要 Newtype 包装？}

\texttt{property.rs} 的属性系统基于 \texttt{TypeId} 区分属性类型：
\[
  \text{属性键} = \texttt{TypeId::of::<T>()}
\]
\texttt{TypeId} 由 Rust 编译器为每个 \texttt{'static} 类型自动分配，相同类型共享同一 ID。
\textbf{问题}：若想为顶点同时附加 \texttt{[f64; 3]} 类型的「法向」与「颜色」，
二者 \texttt{TypeId} 相同，会互相覆盖。

\textbf{解决方案}：用 Newtype 包装，让每个语义属性有独立类型：

\begin{lstlisting}
#[derive(Debug, Clone, Copy, Default, PartialEq)]
pub struct VertexNormal(pub [f64; 3]);

#[derive(Debug, Clone, Copy, Default, PartialEq)]
pub struct VertexColor(pub [f64; 3]);

#[derive(Debug, Clone, Copy, Default, PartialEq)]
pub struct VertexUv(pub [f64; 2]);

#[derive(Debug, Clone, Copy, Default, PartialEq)]
pub struct FaceNormal(pub [f64; 3]);
\end{lstlisting}

\begin{sloppypar}
\texttt{VertexNormal} 与 \texttt{VertexColor} 现在是\textbf{不同类型}，
\texttt{TypeId::of::<VertexNormal>()} $\ne$ \texttt{TypeId::of::<VertexColor>()}，
可以共存于同一 \texttt{MeshProperties} 容器。
\end{sloppypar}

\section{API 表}

\subsection{Newtype 列表}

\begin{center}
\renewcommand{\arraystretch}{1.18}
\begin{tabular}{@{}lll@{}}
  \toprule
  \textbf{Newtype} & \textbf{底层类型} & \textbf{语义} \\
  \midrule
  \texttt{VertexNormal([f64;3])}    & \texttt{[f64;3]} & 顶点法向（单位向量） \\
  \texttt{VertexColor([f64;3])}     & \texttt{[f64;3]} & 顶点颜色 RGB，每通道 $\in [0, 1]$ \\
  \texttt{VertexUv([f64;2])}        & \texttt{[f64;2]} & 顶点 UV 纹理坐标 \\
  \texttt{FaceNormal([f64;3])}      & \texttt{[f64;3]} & 面法向（单位向量） \\
  \texttt{VertexSelected(bool)}     & \texttt{bool}    & 顶点选择态 \\
  \texttt{HalfEdgeSelected(bool)}   & \texttt{bool}    & 半边选择态 \\
  \texttt{FaceSelected(bool)}       & \texttt{bool}    & 面选择态 \\
  \texttt{HalfEdgeColor([f64;3])}   & \texttt{[f64;3]} & 半边/边颜色 RGB \\
  \texttt{FaceColor([f64;3])}       & \texttt{[f64;3]} & 面颜色 RGB \\
  \texttt{VertexSize(f64)}          & \texttt{f64}     & 顶点显示半径 \\
  \texttt{HalfEdgeThickness(f64)}   & \texttt{f64}     & 半边粗细（twin 同步） \\
  \texttt{FaceOpacity(f64)}         & \texttt{f64}     & 面透明度 $\in [0, 1]$ \\
  \bottomrule
\end{tabular}
\end{center}

\subsection{注册函数}

\begin{center}
\renewcommand{\arraystretch}{1.18}
\begin{tabular}{@{}ll@{}}
  \toprule
  \textbf{注册函数} & \textbf{返回类型} \\
  \midrule
  \texttt{add\_vertex\_normal(\&mut props)}    & \texttt{PropHandle<VertexNormal>} \\
  \texttt{add\_vertex\_colors(\&mut props)}    & \texttt{PropHandle<VertexColor>} \\
  \texttt{add\_vertex\_uvs(\&mut props)}       & \texttt{PropHandle<VertexUv>} \\
  \texttt{add\_face\_normal(\&mut props)}      & \texttt{PropHandle<FaceNormal>} \\
  \texttt{add\_vertex\_selection(\&mut props)}  & \texttt{PropHandle<VertexSelected>} \\
  \texttt{add\_halfedge\_selection(\&mut props)} & \texttt{PropHandle<HalfEdgeSelected>} \\
  \texttt{add\_face\_selection(\&mut props)}    & \texttt{PropHandle<FaceSelected>} \\
  \texttt{add\_halfedge\_colors(\&mut props)}   & \texttt{PropHandle<HalfEdgeColor>} \\
  \texttt{add\_face\_colors(\&mut props)}       & \texttt{PropHandle<FaceColor>} \\
  \texttt{add\_vertex\_sizes(\&mut props)}       & \texttt{PropHandle<VertexSize>} \\
  \texttt{add\_halfedge\_thickness(\&mut props)} & \texttt{PropHandle<HalfEdgeThickness>} \\
  \texttt{add\_face\_opacity(\&mut props)}      & \texttt{PropHandle<FaceOpacity>} \\
  \bottomrule
\end{tabular}
\end{center}

\subsection{批量计算与收集}

\begin{center}
\renewcommand{\arraystretch}{1.18}
\begin{tabular}{@{}ll@{}}
  \toprule
  \textbf{函数} & \textbf{语义} \\
  \midrule
  \texttt{populate\_vertex\_normal(...)} & 用 \texttt{vertex\_normal} 填充所有顶点 \\
  \texttt{populate\_face\_normal(...)}   & 用 \texttt{face\_normal} 填充所有面 \\
  \texttt{collect\_vertex\_normal(...)}  & 按顶点顺序提取法向 \\
  \texttt{collect\_vertex\_uvs(...)}     & 按顶点顺序提取 UV \\
  \texttt{collect\_vertex\_colors(...)}  & 按顶点顺序提取颜色 \\
  \bottomrule
\end{tabular}
\end{center}

\subsection{OBJ 属性感知}

\begin{center}
\renewcommand{\arraystretch}{1.18}
\begin{tabular}{@{}ll@{}}
  \toprule
  \textbf{函数} & \textbf{语义} \\
  \midrule
  \texttt{parse\_obj\_with\_attrs(text)}   & 解析 OBJ 并提取 \texttt{vt}/\texttt{vn} \\
  \texttt{format\_obj\_with\_attrs(...)}   & 序列化带 \texttt{vt}/\texttt{vn} 的 OBJ \\
  \texttt{install\_vertex\_attrs(...)}     & 一次性注册并填充 \\
  \bottomrule
\end{tabular}
\end{center}

\section{属性数据流}

\begin{center}
\begin{tikzpicture}[
  node distance=0.7cm,
  proc/.style={draw, rounded corners=2pt, minimum width=4.5cm, minimum height=1.2cm, align=center, font=\small},
  op/.style={-Stealth, thick},
  store/.style={draw, cylinder, shape border rotate=90, aspect=0.3, minimum height=1.5cm, minimum width=3cm, font=\small, fill=blue!8}
]
  \node[store] (obj) {OBJ 文本};
  \node[proc, right=1.6cm of obj, fill=yellow!15] (parser) {\texttt{\scriptsize parse\_obj\_with\_attrs} \\ 解析 v / vt / vn / f};
  \node[proc, right=1.6cm of parser, fill=green!12] (mesh) {\texttt{MeshStorage} \\ + \texttt{MeshProperties}};
  \node[proc, below=1.2cm of mesh, fill=red!10] (writer) {\texttt{\scriptsize format\_obj\_with\_attrs} \\ 写出 v / vt / vn / f};
  \node[store, left=1.6cm of writer] (out) {OBJ 文本};

  \draw[op] (obj) -- (parser);
  \draw[op] (parser) -- (mesh);
  \draw[op] (mesh) -- (writer);
  \draw[op] (writer) -- (out);
\end{tikzpicture}
\end{center}

\section{OBJ 属性感知 IO}

\subsection{解析：parse\_obj\_with\_attrs}

输入：OBJ 文本，含 \texttt{v} / \texttt{vt} / \texttt{vn} / \texttt{f} 行。
输出：\texttt{(MeshStorage, MeshProperties)}，属性系统已注册
\texttt{VertexNormal} 与 \texttt{VertexUv}（仅在输入含对应行时）。

\paragraph{面行格式} 支持 4 种：
\begin{itemize}
  \item \texttt{f v1 v2 v3} —— 仅顶点；
  \item \texttt{f v1/vt1 v2/vt2 v3/vt3} —— 顶点 + UV；
  \item \texttt{f v1/vt1/vn1 v2/vt2/vn2 v3/vt3/vn3} —— 顶点 + UV + 法向；
  \item \texttt{f v1//vn1 v2//vn2 v3//vn3} —— 顶点 + 法向（无 UV）。
\end{itemize}

\paragraph{属性映射} 每个面顶点 \texttt{(vt, vn)} 索引（1-based）映射到顶点：
\begin{enumerate}
  \item 收集所有 \texttt{vt} 行到 \texttt{texcoords: Vec<[f64;2]>}；
  \item 收集所有 \texttt{vn} 行到 \texttt{normals: Vec<[f64;3]>}；
  \item 遍历每个面顶点 {\scriptsize\texttt{(v\_idx, vt\_idx, vn\_idx)}}：
        {\scriptsize\texttt{vert\_normal[v\_idx] = normals[vn\_idx - 1]}}；
  \item 按顶点顺序注册属性并填充 \texttt{MeshProperties}。
\end{enumerate}

\paragraph{同顶点多面引用} 同一顶点可能被多个面引用，
\textbf{后面的面覆盖前面的}（OBJ 标准允许同顶点在不同面有不同 UV/法向，
本实现按顶点而非按面顶点存储属性，故取最后一次赋值）。

\subsection{序列化：format\_obj\_with\_attrs}

输出顺序：
\begin{enumerate}
  \item 所有 \texttt{v} 行（顶点位置）；
  \item 所有 \texttt{vt} 行（若 \texttt{uv\_handle} 为 \texttt{Some}）；
  \item 所有 \texttt{vn} 行（若 \texttt{normal\_handle} 为 \texttt{Some}）；
  \item 所有 \texttt{f} 行，按 \texttt{v/vt/vn} 格式输出（依据已注册属性选择）。
\end{enumerate}

\paragraph{面行格式选择}
\begin{itemize}
  \item 同时有 UV 和法向：\texttt{f v/vt/vn v/vt/vn v/vt/vn}；
  \item 仅有 UV：\texttt{f v/vt v/vt v/vt}；
  \item 仅有法向：\texttt{f v//vn v//vn v//vn}；
  \item 都无：\texttt{f v v v}。
\end{itemize}

未设置属性的顶点写 0 占位（\texttt{vt 0 0} 或 \texttt{vn 0 0 0}）。

\section{批量安装：install\_vertex\_attrs}

\texttt{install\_vertex\_attrs(props, normals, uvs, colors, vertex\_ids)}
用于从外部数据（如其他 IO 库解析结果）一次性注册并填充属性：

\begin{center}
\begin{tikzpicture}[
  node distance=0.55cm,
  proc/.style={draw, rounded corners=2pt, minimum width=10cm, minimum height=0.7cm, align=center, font=\small},
  op/.style={-Stealth, thick},
  startend/.style={draw, rounded corners=10pt, minimum width=2.6cm, minimum height=0.7cm, font=\small, fill=gray!12}
]
  \node[startend] (s) {开始};
  \node[proc, below=of s, fill=blue!8] (p1) {若 normals 为 \texttt{Some}：注册 \texttt{VertexNormal}，按索引填充};
  \node[proc, below=of p1, fill=blue!8] (p2) {若 uvs 为 \texttt{Some}：注册 \texttt{VertexUv}，按索引填充};
  \node[proc, below=of p2, fill=blue!8] (p3) {若 colors 为 \texttt{Some}：注册 \texttt{VertexColor}，按索引填充};
  \node[startend, below=of p3] (e) {返回};
  \draw[op] (s) -- (p1);
  \draw[op] (p1) -- (p2);
  \draw[op] (p2) -- (p3);
  \draw[op] (p3) -- (e);
\end{tikzpicture}
\end{center}

\textbf{部分输入}：任一参数为 \texttt{None} 则跳过对应属性，
\texttt{has\_vertex\_prop::<T>()} 对应返回 \texttt{false}。

\section{与 property.rs 的关系}

\begin{center}
\renewcommand{\arraystretch}{1.18}
\begin{tabular}{@{}lll@{}}
  \toprule
  \textbf{维度} & \texttt{property.rs}（泛型） & \texttt{builtin\_attrs.rs}（内置） \\
  \midrule
  类型      & 任意 \texttt{T: 'static}              & 固定 4 种 Newtype \\
  句柄      & \texttt{PropertyHandle<T>}            & 同左（直接复用） \\
  存储      & \texttt{HashMap<TypeId, Box<dyn ErasedProperty>>} & 同左 \\
  注册      & \texttt{add\_vertex\_prop::<T>()}     & \texttt{add\_vertex\_normal(\&mut props)} \\
  读写      & \texttt{get/set\_vertex\_prop}        & 同左 \\
  语义      & 无（用户自定义）                       & 法向 / UV / 颜色 / 面法向 \\
  IO 联动   & 无                                    & \texttt{parse\_obj\_with\_attrs} 等 \\
  \bottomrule
\end{tabular}
\end{center}

\texttt{builtin\_attrs} 完全建立在 \texttt{property.rs} 之上，不修改其底层存储。
用户可同时使用两者：自定义泛型属性与内置属性互不冲突。

\section{退化守卫}

\begin{center}
\renewcommand{\arraystretch}{1.18}
\begin{tabular}{@{}lll@{}}
  \toprule
  \textbf{函数} & \textbf{退化场景} & \textbf{处理} \\
  \midrule
  \texttt{populate\_vertex\_normal} & 顶点法向不可计算（孤立 / 边界退化） & 跳过该顶点 \\
  \texttt{populate\_face\_normal}   & 面法向不可计算（共线） & 跳过该面 \\
  \texttt{populate\_face\_normal}   & 空网格 & 安全返回（无 panic） \\
  \texttt{collect\_vertex\_normal}  & 顶点未设置法向 & 填 \texttt{[0, 0, 0]} \\
  \texttt{collect\_vertex\_colors}  & 顶点未设置颜色 & 填 \texttt{[1, 1, 1]}（白） \\
  \texttt{collect\_vertex\_uvs}     & 顶点未设置 UV & 填 \texttt{[0, 0]} \\
  \texttt{parse\_obj\_with\_attrs}  & OBJ 无 \texttt{vn} 行 & 不注册 \texttt{VertexNormal} \\
  \texttt{parse\_obj\_with\_attrs}  & 面索引越界 & 返回 \texttt{IndexOutOfRange} 错误 \\
  \texttt{parse\_obj\_with\_attrs}  & 面顶点数 $< 3$ & 返回 \texttt{NotTriangular} 错误 \\
  \bottomrule
\end{tabular}
\end{center}

\section{使用示例}

\subsection{为 icosphere 计算顶点法向并写出 OBJ}

\begin{lstlisting}
use halfedge::*;
use halfedge::builtin_attrs::*;

let mesh = test_util::build_icosphere(2);
let mut props = MeshProperties::new();
let n_handle = add_vertex_normals(&mut props);
populate_vertex_normals(&mesh, &mut props, n_handle);

let obj = format_obj_with_attrs(&mesh, &props,
    Some(n_handle), None);
// obj 现在含 v / vn / f 行，f 行格式为 "f v//vn v//vn v//vn"
\end{lstlisting}

\subsection{从带属性的 OBJ 加载}

\begin{lstlisting}
let text = std::fs::read_to_string("model.obj")?;
let (mesh, props) = parse_obj_with_attrs(&text)?;
// props 已注册 VertexNormal 与 VertexUv（若 OBJ 含 vn / vt 行）
\end{lstlisting}

\subsection{批量安装自定义属性}

\begin{lstlisting}
let mesh = build_icosphere(0);
let ids: Vec<VertexId> = mesh.vertex_ids().collect();
let mut props = MeshProperties::new();

let normals: Vec<[f64; 3]> = ids.iter().map(|_| [0.0, 1.0, 0.0]).collect();
let uvs: Vec<[f64; 2]> = ids.iter().enumerate()
    .map(|(i, _)| [i as f64 * 0.1, 0.0])
    .collect();

install_vertex_attrs(&mut props, Some(&normals), Some(&uvs), None, &ids);
assert!(props.has_vertex_prop::<VertexNormal>());
assert!(props.has_vertex_prop::<VertexUv>());
assert!(!props.has_vertex_prop::<VertexColor>());
\end{lstlisting}

\section{选择态与边/面色}
\label{sec:selection-color}

本节描述为顶点 / 半边 / 面三类元素提供的「选择态」与「边/面颜色」便捷 API。
所有 API 都建立在 \texttt{MeshProperties} 之上，复用 Newtype 模式：

\subsection{Newtype 列表}

\begin{center}
\renewcommand{\arraystretch}{1.18}
\begin{tabular}{@{}lll@{}}
  \toprule
  \textbf{Newtype} & \textbf{底层类型} & \textbf{语义} \\
  \midrule
  \texttt{VertexSelected}   & \texttt{bool}    & 顶点是否被选中 \\
  \texttt{HalfEdgeSelected} & \texttt{bool}    & 半边是否被选中 \\
  \texttt{FaceSelected}     & \texttt{bool}    & 面是否被选中 \\
  \texttt{HalfEdgeColor}    & \texttt{[f64;3]} & 半边/边颜色 RGB \\
  \texttt{FaceColor}        & \texttt{[f64;3]} & 面颜色 RGB \\
  \bottomrule
\end{tabular}
\end{center}

\subsection{语义约定}

\begin{itemize}
  \item \textbf{未注册 $\equiv$ false/None}：未注册属性时，\texttt{is\_*\_selected}
        返回 \texttt{false}，\texttt{edge\_color}/\texttt{face\_color} 返回 \texttt{None}，
        避免强制初始化。
  \item \textbf{按类型清理}：\texttt{clear\_*\_selection} 与 \texttt{clear\_face\_colors}
        仅删除当前 Newtype 类型的属性，\textbf{不影响其他类型}（如 \texttt{VertexNormal}）。
  \item \textbf{懒迭代}：\texttt{selected\_*\_ids} 返回的迭代器按需过滤，
        不预收集。
\end{itemize}

\subsection{边级（twin 对）同步}

\begin{sloppypar}
无向边由两条互为 \texttt{twin} 的半边组成。从语义上，边的选择/颜色应一致。
\texttt{select\_edge} / \texttt{deselect\_edge} / \texttt{set\_edge\_color}
\textbf{同时写入} \texttt{he} 与其 \texttt{twin}；\texttt{is\_edge\_selected}
\textbf{任一}半边被选中即返回 \texttt{true}；\texttt{selected\_edge\_ids}
按 canonical half（\texttt{he $\le$ twin}）去重，每个无向边只返回一次。
\end{sloppypar}

\begin{center}
\begin{tikzpicture}[
  node distance=1.0cm,
  vert/.style={circle, draw=blue!70!black, fill=blue!10, minimum size=7mm, inner sep=0pt, font=\small},
  he/.style={-Stealth, thick, draw=red!70!black},
  lab/.style={font=\footnotesize, sloped, above}
]
  \node[vert] (v0) {$v_0$};
  \node[vert, right=3.2cm of v0] (v1) {$v_1$};
  \draw[he] (v0) -- node[lab] {$h$（选中）} (v1);
  \draw[he, bend left=35] (v1) -- node[lab, below] {$h.\mathrm{twin}$（自动同步）} (v0);
\end{tikzpicture}
\end{center}

\subsection{API 表}

\begin{longtable}{@{}ll@{}}
  \toprule
  \textbf{API} & \textbf{语义} \\
  \midrule
  \endhead
  \multicolumn{2}{l}{\textit{顶点选择}} \\
  \texttt{select\_vertex(props, h, id)}                       & 选中顶点 \\
  \texttt{deselect\_vertex(props, h, id)}                    & 取消选中 \\
  \texttt{toggle\_vertex\_selection(props, h, id) $\to$ bool} & 切换并返回新状态 \\
  \texttt{is\_vertex\_selected(props, h, id) $\to$ bool}     & 查询 \\
  \texttt{clear\_vertex\_selection(mesh, props, h)}          & 清空所有顶点选择 \\
  \texttt{selected\_vertex\_ids(mesh, props, h) $\to$ Iterator} & 懒迭代选中顶点 \\
  \texttt{select\_all\_vertices(mesh, props, h)}             & 全选 \\
  \texttt{invert\_vertex\_selection(mesh, props, h)}        & 反选 \\
  \texttt{count\_selected\_vertices(mesh, props, h) $\to$ usize} & 统计 \\
  \midrule
  \multicolumn{2}{l}{\textit{半边选择}} \\
  \texttt{select\_halfedge / deselect\_halfedge}             & 同半边 \\
  \texttt{toggle\_halfedge\_selection}                       & 同半边 \\
  \texttt{is\_halfedge\_selected}                            & 同半边 \\
  \texttt{clear\_halfedge\_selection(mesh, props, h)}        & 清空 \\
  \texttt{selected\_halfedge\_ids(mesh, props, h) $\to$ Iterator} & 懒迭代 \\
  \midrule
  \multicolumn{2}{l}{\textit{面选择}} \\
  \texttt{select\_face / deselect\_face}                     & 同面 \\
  \texttt{toggle\_face\_selection}                           & 同面 \\
  \texttt{is\_face\_selected}                                & 同面 \\
  \texttt{clear\_face\_selection(mesh, props, h)}           & 清空 \\
  \texttt{selected\_face\_ids(mesh, props, h) $\to$ Iterator} & 懒迭代 \\
  \texttt{select\_all\_faces(mesh, props, h)}                & 全选 \\
  \texttt{count\_selected\_faces(mesh, props, h) $\to$ usize} & 统计 \\
  \midrule
  \multicolumn{2}{l}{\textit{边级（twin 同步）}} \\
  \texttt{select\_edge(mesh, props, h, he)}                  & 选中 \texttt{he} 与 \texttt{twin} \\
  \texttt{deselect\_edge(mesh, props, h, he)}                & 取消 \texttt{he} 与 \texttt{twin} \\
  \texttt{toggle\_edge\_selection(mesh, props, h, he) $\to$ bool} & 切换 \\
  \texttt{is\_edge\_selected(mesh, props, h, he) $\to$ bool} & 任一 half 选中即 true \\
  \texttt{selected\_edge\_ids(mesh, props, h) $\to$ Iterator} & canonical 去重 \\
  \texttt{set\_edge\_color(mesh, props, h, he, color)}       & 同步设置 \texttt{he} 与 \texttt{twin} 颜色 \\
  \texttt{edge\_color(mesh, props, h, he) $\to$ Option<[f64;3]>} & 优先 \texttt{he}，未设置则尝试 \texttt{twin} \\
  \midrule
  \multicolumn{2}{l}{\textit{面颜色}} \\
  \texttt{set\_face\_color(props, h, id, color)}              & 设置面颜色 \\
  \texttt{face\_color(props, h, id) $\to$ Option<[f64;3]>}    & 查询 \\
  \texttt{clear\_face\_colors(mesh, props, h)}                & 清空 \\
  \bottomrule
\end{longtable}

\subsection{使用示例}

\begin{lstlisting}
use halfedge::*;
let mesh = build_icosphere(1);
let mut props = MeshProperties::new();

// 选中所有边界顶点
let hv = add_vertex_selection(&mut props);
for v in mesh.vertex_ids() {
    if is_boundary_vertex(&mesh, v) {
        select_vertex(&mut props, hv, v);
    }
}
let n_boundary = count_selected_vertices(&mesh, &props, hv);

// 选中所有内部边（仅示例：取代表性边）
let he = add_halfedge_selection(&mut props);
let some_edge = mesh.halfedge_ids().next().unwrap();
select_edge(&mesh, &mut props, he, some_edge);
assert!(is_edge_selected(&mesh, &props, he, some_edge));

// 为每个面着色
let fc = add_face_colors(&mut props);
for (i, f) in mesh.face_ids().enumerate() {
    let c = if i % 2 == 0 { [0.8, 0.2, 0.2] } else { [0.2, 0.8, 0.2] };
    set_face_color(&mut props, fc, f, c);
}
\end{lstlisting}

\section{可视化尺寸 / 粗细 / 透明度}
\label{sec:size-thickness-opacity}

本节描述为顶点 / 半边 / 面三类元素提供的「可视化尺寸」属性，用于控制渲染时
的几何外观。所有 API 与颜色 API 形态一致，建立在 \texttt{MeshProperties} 之上。

\subsection{Newtype 列表}

\begin{longtable}{@{}llp{5.2cm}@{}}
  \toprule
  \textbf{Newtype} & \textbf{底层类型} & \textbf{语义} \\
  \midrule
  \endhead
  \texttt{VertexSize}        & \texttt{f64} & 顶点显示半径（单位与 \texttt{position} 一致），默认 0 = 不渲染标记 \\
  \texttt{HalfEdgeThickness} & \texttt{f64} & 半边/边线宽（twin 同步） \\
  \texttt{FaceOpacity}       & \texttt{f64} & 面透明度 $\in [0, 1]$：0 = 完全透明，1 = 完全不透明 \\
  \bottomrule
\end{longtable}

\subsection{语义约定}

\begin{sloppypar}
\begin{itemize}
  \item \textbf{未注册 $\equiv$ None}：未注册属性时，\texttt{vertex\_size} / \texttt{edge\_thickness} / \texttt{face\_opacity} 返回 \texttt{None}，避免强制初始化。
  \item \textbf{clamp 透明度}：\texttt{set\_face\_opacity} 与 \texttt{set\_uniform\_face\_opacity} 自动将输入 clamp 到 $[0, 1]$，杜绝越界。
  \item \textbf{按类型清理}：\texttt{clear\_*} 仅删除当前 Newtype 类型的属性，不影响其他类型。
  \item \textbf{批量统一}：\texttt{set\_uniform\_*} 一行代码为所有元素设置同一值。
\end{itemize}
\end{sloppypar}

\subsection{边级（twin 对）粗细同步}

与边颜色 API 一致，\texttt{set\_edge\_thickness} 同时写入 \texttt{he} 与其 \texttt{twin}；
\texttt{edge\_thickness} 优先读 \texttt{he}，未设置则 fallback 到 \texttt{twin}；
\texttt{set\_uniform\_edge\_thickness} 仅处理 canonical half（\texttt{he $\le$ twin}），
避免重复设置。

\begin{sloppypar}
若需非对称渲染（两条半边不同粗细），可使用 \texttt{set\_halfedge\_thickness} /
\texttt{halfedge\_thickness} 单半边版本。
\end{sloppypar}

\subsection{API 表}

\begin{longtable}{@{}ll@{}}
  \toprule
  \textbf{API} & \textbf{语义} \\
  \midrule
  \endhead
  \multicolumn{2}{l}{\textit{顶点大小}} \\
  \texttt{set\_vertex\_size(props, h, id, size)}                & 设置顶点显示半径 \\
  \texttt{vertex\_size(props, h, id) $\to$ Option<f64>}        & 查询 \\
  \texttt{clear\_vertex\_sizes(mesh, props, h)}                & 清空 \\
  \texttt{set\_uniform\_vertex\_size(mesh, props, h, size)}    & 批量统一 \\
  \midrule
  \multicolumn{2}{l}{\textit{边粗细（twin 同步）}} \\
  \texttt{set\_edge\_thickness(mesh, props, h, he, t)}          & 同步设置 \texttt{he} 与 \texttt{twin} \\
  \texttt{edge\_thickness(mesh, props, h, he) $\to$ Option<f64>} & 优先 \texttt{he}，未设置则尝试 \texttt{twin} \\
  \texttt{set\_halfedge\_thickness(props, h, he, t)}             & 单半边设置（不同步 twin） \\
  \texttt{halfedge\_thickness(props, h, he) $\to$ Option<f64>}  & 单半边查询 \\
  \texttt{clear\_halfedge\_thickness(mesh, props, h)}           & 清空 \\
  \texttt{set\_uniform\_edge\_thickness(mesh, props, h, t)}    & 批量统一（canonical 去重） \\
  \midrule
  \multicolumn{2}{l}{\textit{面透明度}} \\
  \texttt{set\_face\_opacity(props, h, id, op)}                & 设置面透明度（自动 clamp） \\
  \texttt{face\_opacity(props, h, id) $\to$ Option<f64>}       & 查询 \\
  \texttt{clear\_face\_opacity(mesh, props, h)}                & 清空 \\
  \texttt{set\_uniform\_face\_opacity(mesh, props, h, op)}     & 批量统一（自动 clamp） \\
  \bottomrule
\end{longtable}

\subsection{使用示例}

\begin{lstlisting}
use halfedge::*;
let mesh = build_icosphere(1);
let mut props = MeshProperties::new();

// 统一设置顶点大小、边粗细、面透明度
let hv = add_vertex_sizes(&mut props);
set_uniform_vertex_size(&mesh, &mut props, hv, 0.02);

let he = add_halfedge_thickness(&mut props);
set_uniform_edge_thickness(&mesh, &mut props, he, 0.005);

let hf = add_face_opacity(&mut props);
set_uniform_face_opacity(&mesh, &mut props, hf, 0.85);

// 单独调整一个面为半透明高亮
let f0 = mesh.face_ids().next().unwrap();
set_face_opacity(&mut props, hf, f0, 0.5);
assert_eq!(face_opacity(&props, hf, f0), Some(0.5));

// 越界输入自动 clamp
set_face_opacity(&mut props, hf, f0, 1.5);
assert_eq!(face_opacity(&props, hf, f0), Some(1.0));
\end{lstlisting}

\section{测试覆盖}

\begin{longtable}{@{}ll@{}}
  \toprule
  \textbf{测试名} & \textbf{验证点} \\
  \midrule
  \endhead
  \multicolumn{2}{l}{\textit{基础属性}} \\
  \texttt{vertex\_normal\_roundtrip}        & icosphere 顶点法向为单位向量 \\
  \texttt{face\_normal\_roundtrip}          & 所有面法向已写入 \\
  \texttt{vertex\_color\_default\_is\_white} & 未设置颜色返回默认白 \\
  \texttt{vertex\_uv\_set\_and\_get}        & UV 设置后可读取 \\
  \texttt{install\_vertex\_attrs\_handles\_partial\_input} & 仅传法向不注册 UV/Color \\
  \texttt{install\_vertex\_attrs\_all\_three} & 三种属性同时安装 \\
  \texttt{newtype\_default\_is\_zero}       & Newtype 默认值为零 \\
  \texttt{missing\_handle\_returns\_none}   & 未注册句柄读取返回 \texttt{None} \\
  \texttt{populate\_face\_normal\_skips\_invalid} & 空网格不 panic \\
  \midrule
  \multicolumn{2}{l}{\textit{选择态与边/面色}} \\
  \texttt{vertex\_selection\_basic}                     & select / toggle / is / count \\
  \texttt{vertex\_select\_all\_and\_invert}             & 全选 + 反选 \\
  \texttt{vertex\_clear\_selection\_keeps\_registration} & 清空后类型注册保留 \\
  \texttt{selected\_vertex\_ids\_iterator\_correct}     & 懒迭代器返回选中顶点 \\
  \texttt{face\_selection\_basic}                      & 面 select / select\_all / clear \\
  \texttt{halfedge\_selection\_basic}                   & 半边 select / toggle / clear \\
  \texttt{edge\_selection\_syncs\_twin\_pair}           & 选中半边时 twin 同步 \\
  \texttt{selected\_edge\_ids\_returns\_canonical\_only} & canonical half 去重 \\
  \texttt{edge\_color\_syncs\_twin\_pair}               & 颜色同步到 twin \\
  \texttt{face\_color\_basic}                           & 面颜色 set/get/clear \\
  \texttt{selection\_newtype\_default\_is\_false}       & 选择态默认 false \\
  \texttt{color\_newtype\_default\_is\_zero}            & 颜色默认 $[0,0,0]$ \\
  \texttt{unregistered\_selection\_safe\_no\_panic}    & 未注册查询不 panic \\
  \texttt{unregistered\_edge\_safe\_no\_panic}          & 未注册边查询不 panic \\
  \texttt{empty\_mesh\_selection\_no\_panic}           & 空网格调用全选/反选不 panic \\
  \midrule
  \multicolumn{2}{l}{\textit{可视化尺寸 / 粗细 / 透明度}} \\
  \texttt{vertex\_size\_basic}                              & set / uniform / clear + 类型注册保留 \\
  \texttt{vertex\_size\_default\_is\_zero}                  & Newtype 默认值为 0 \\
  \texttt{edge\_thickness\_syncs\_twin\_pair}              & twin 同步 + fallback 读取 \\
  \texttt{edge\_thickness\_uniform\_sets\_all}             & 批量统一设置所有边 \\
  \texttt{halfedge\_thickness\_independent\_set}          & 单半边设置不影响 twin \\
  \texttt{face\_opacity\_basic}                            & set / clamp / uniform / clear \\
  \texttt{face\_opacity\_default\_is\_zero}                & Newtype 默认值为 0 \\
  \texttt{unregistered\_size\_thickness\_opacity\_safe} & 未注册查询不 panic \\
  \texttt{empty\_mesh\_size\_thickness\_opacity} & 空网格批量操作不 panic \\
  \bottomrule
\end{longtable}

\texttt{src/builtin\_attrs.rs} 共 \textbf{33 个}单元测试，全库共 \textbf{708 个}。

\end{document}
